% ===========================================================================
\section{Introduction}
\label{sec:intro}
% ===========================================================================

A national greenhouse-gas inventory is a large arithmetic object whose correctness depends
almost entirely on properties that span its parts. Fuel supplied to a country must be
allocated among the sectors that burn it without any of it being counted twice or dropped.
Nitrogen removed from one balance must arrive in another. A category reported under one IPCC
code must not also appear under a second. Each of these is a statement about the relationship
between two sectors, and in the conventional production model each is enforced by a compiler
reading two spreadsheets side by side and satisfying themselves that the numbers agree.

That model has a failure mode that is difficult to detect from the outside and easy to
introduce from the inside: a change to one sector silently invalidates a property that was
true when a different sector was written. Nothing breaks. The total still sums. The
inconsistency is discoverable only by repeating the original reading, which nobody does,
because the reading was tedious the first time and there is no signal that it needs repeating.

\subsection{The argument of this paper}

We built Pakistan's inventory so that those cross-cutting properties are executable. Each
sector is a Python module satisfying a common contract; the contract is enforced by
\num{32} tests that run against every registered sector; the properties that span sectors
are assertions with \texttt{==} in the sector code and in the test suite, not sentences in a
methods document. When a change to the agriculture sector's manure-management shares alters
how much dung is diverted from the manure balance as fuel, the residential sector's dung ledger
changes with it, because it is derived from the same herd through the same volatile-solids
rates rather than estimated separately. The two cannot drift apart, because there is only one
number.

This is a software-engineering claim about inventory compilation, and it is testable in the
ordinary way: the repository is public, the test suite runs in continuous integration, and
every figure in this paper is regenerated from the engine's own output files by a build step
that fails if a figure is stale.

\subsection{Why Pakistan, and why now}

Pakistan is the fifth most populous country in the world and among the most exposed to
climate impacts, yet its inventory infrastructure is thin. Four official filings are in
circulation and they are not interchangeable, so this paper distinguishes them throughout.
The First Biennial Update Report \citep{pakbur1} carries the most recent full narrative
inventory with a published sectoral breakdown, for FY2017--18, and is the comparator for
categories where only a narrative figure exists. The First Biennial Transparency Report
\citep{pakbtr2024} supersedes it as the current submission and carries Pakistan's first
Common Reporting Tables \citep{pakcrt2026}, for inventory year 2022; those tables are
machine-readable at cell level and are the comparator used wherever a category cell exists,
because a cell can be audited and a narrative total cannot. The Third National Communication
\citep{paknc32025} and the Initial National Communication complete the set; the latter is
used once, as a labelled magnitude check in the railway benchmark register
(Section~\ref{sec:benchmarks}), and nowhere else. Where this paper says a category is compared against
Pakistan's own reporting, it names which of these it means, and Section~\ref{sec:benchmarks}
prefers the CRT cells over the BUR-1 narrative wherever both exist.

That distinction is not bookkeeping. The CRT's own notation keys record that residential
biomass---dung cake and firewood, the country's two largest household fuels---is NE in the
national filing, so a comparison against its residential total is a comparison against the
metered-fossil sliver of the sector rather than against the sector.

International datasets fill the remaining gap---EDGAR \citep{edgar2025booklet}, FAOSTAT
\citep{faostat}, Climate TRACE \citep{climatetrace}---but each is a global product
applying globally uniform assumptions to nationally reported activity, and their Pakistani
sector estimates disagree with each other by amounts that are large relative to the
mitigation decisions they inform.

There is also a specific structural reason a national inventory for Pakistan is worth
building as a modular system rather than a document. Its emissions are unusually
concentrated in agriculture---enteric fermentation alone is \SI{34.1}{\percent} of the
national total, larger than every energy category combined except road transport and
electricity---which means an inventory dominated by one sector must nevertheless get the
other nine right to be useful for anything except a headline. Sectors that are small in
\COe{} can be large in policy: residential biomass is \SI{5.4}{\percent} of national
emissions and the largest estimated source of gross primary \PM{} emissions
among the sectors quantified here. Whether it is also the dominant source of \emph{exposure}
is a separate question this inventory cannot answer: exposure depends on ventilation,
proximity and time-activity patterns that are outside its scope.

\subsection{Contributions}

\begin{enumerate}[leftmargin=*,itemsep=2pt,topsep=3pt]
\item \textbf{A ten-module national inventory covering the principal currently
  reconstructable anthropogenic sources} for Pakistan at
  \SI{444.96}{\mega\tonne} \COe{} under reference scenario \texttt{REF\_2023\_24}---\SI{87}{\percent} of the
  corresponding EDGAR aggregate magnitude, with every mapped IPCC category estimated over
  \SI{11}{\percent} of that scope and the unestimated categories named over the rest---in
  which every
  input carries an explicit provenance and evidence-status record, Pakistani official sources
  are used where available and proxy, survey-derived or assumed values are identified
  separately, every sector is allocated to provinces, and IPCC notation keys are reported
  over all \num{138} category--gas combinations considered.
\item \textbf{A sector-as-module architecture} in which \num{32} contract invariants are
  enforced on every sector, so that adding a sector cannot silently weaken the inventory
  (Section~\ref{sec:architecture}).
\item \textbf{Four enforced accounting closures} asserted by construction rather than
  described, each of which corrects a double-count or an orphaned flow that the prose version
  of the same inventory had already documented as a known defect. \emph{Two of the four cross
  a module}---dung nitrogen from agriculture into residential, and generator petrol carved out
  of transport's control by residential---and two are between categories inside one module,
  which is why they are called accounting closures here rather than cross-sector ones
  (Section~\ref{sec:closures}).
\item \textbf{A methodological finding about the IPCC default tables}: four instances in
  which the plausible wrong table leaves \COtwo{} exactly correct and is therefore invisible
  to every external comparator, moving the affected category by \SI{8}{\percent} where the
  category is small enough for a trace gas to matter and by less than
  \SI{0.5}{\percent} where it is \COtwo{}-dominated (Section~\ref{sec:factors}).
\item \textbf{Uncertainty carried at the level of parameters rather than of emission rows},
  so that one row may depend on an activity, a conversion and a factor at once, each shared
  with a different set of other rows (Section~\ref{sec:uncertainty}). \emph{This overturns a
  finding of the previous edition.} Under the earlier model, in which a row carried a single
  correlation group, sharing the fuel-carbon factors widened the road-transport category by
  more than half; at parameter level it does not move it at all, because the carbon content of
  diesel carries about \SI{2}{\percent} and the road rows carry thirty to sixty, and
  correlating a small part of a large variance cannot widen the sum. What the parameter graph now
  finds is that \SI{2.5}{\percent} of the national variance has no named parameter, after
  agriculture's IPCC factor
  uncertainties (composed inside each row's declared width, invisible to the registry) were
  given names and the remaining sectors' factor,
  conversion and activity groups were scoped to their source streams the same way, with no
  central number changed in either step. Where the variance sits is read from the released
  ensemble: enteric fermentation carries \SI{59.8}{\percent} of it. Parameter-level
  importance is not claimed.
\item \textbf{A national key-category analysis ranked by variance as well as by level},
  identifying where measurement effort reduces national uncertainty rather than where the
  emissions are (Section~\ref{sec:keycat}).
\end{enumerate}

The agriculture sector's methods, its geospatial allocation and its air-pollutant layer are
documented separately \citep{ahmad2026ag}, as is the nowcasting system that extends it beyond
the census year \citep{ahmad2026nowcast}. This paper treats agriculture as one of ten
sectors and does not repeat that material. Eight of the nine non-agricultural sectors also
have a full technical report---data provenance, methods, factors, validation, uncertainty,
air pollutants and limitations---maintained in the repository, compiled on every release
and held to the same generated-from-the-engine discipline as this manuscript; industrial
fuel combustion was added after those reports were written and is documented in
Section~\ref{sec:completeness} and the completeness table rather than in a standalone
report. Those reports
are versioned documentation of the release rather than companion submissions: this is the
one article, and it carries what a reader needs in order to judge the dataset, while they
carry what a reader needs in order to audit and maintain it.
